% !TEX root = ../../main.tex
\subsection{Exercise 3}
\task{Requirement}
To allocate memory, please refer to the following syscall:
\begin{lstlisting}
    li $v0, 9   # system call code for dynamic allocation
    li $a0, 24  # $a0 contains number of bytes to allocate
\end{lstlisting}
After the above system call, \texttt{\$v0} contains the first address in heap memory
that is allocated. Then, accessing the allocated memory can be done by
\texttt{lw}/\texttt{sw}, for example:
\begin{lstlisting}
    # Trying to write to allocated space
    addi $t0, $zero, 2021
    sw $t0, 0($v0)
\end{lstlisting}
The followings are instructions used to access a file (open/close/read/write):
\begin{lstlisting}
    # Sample MIPS program that writes to a new file.
    # by Kenneth Vollmar and Pete Sanderson
    .data
    fout: .asciiz "testout.txt" # filename for output
    msg1: .asciiz "Before read: "
    msg2: .asciiz "After read: "
    buffer_write: .asciiz "The quick brown fox jumps over the lazy dog.\n"
    buffer_read: .asciiz "---------------------------------\n"
    .text
    #######################################################
    # Open (for writing) a file that does not exist
    li $v0, 13      # system call for open file
    la $a0, fout    # output file name
    li $a1, 1       # Open for writing (flags are 0: read, 1: write)
    li $a2, 0       # mode is ignored
    syscall         # open a file (file descriptor returned in $v0)
    move $s6, $v0   # save the file descriptor
    #######################################################
    # Write to file just opened
    li $v0, 15      # system call for write to file
    move $a0, $s6   # file descriptor
    la $a1, buffer_write # address of buffer from which to write
    li $a2, 44      # hardcoded buffer length
    syscall         # write to file
    #######################################################
    # Close the file
    li $v0, 16      # system call for close file
    move $a0, $s6   # file descriptor to close
    syscall         # close file
    #######################################################
    #######################################################
    # Open (for reading) a file
    li $v0, 13      # system call for open file
    la $a0, fout    # input file name
    li $a1, 0       # Open for reading (flags are 0: read, 1: write)
    li $a2, 0       # mode is ignored
    syscall         # open a file (file descriptor returned in $v0)
    move $s6, $v0   # save the file descriptor
    #######################################################
    # Read from file
    li $v0, 14      # system call for read
    move $a0, $s6   # file descriptor
    la $a1, buffer_read # address of buffer read
    li $a2, 44      # hardcoded buffer length
    syscall         # read file
\end{lstlisting}
Please do the followings:
\begin{sloppypar}
\begin{enumerate}
  \item Manually create a line of text in format
        \texttt{<name>,<id>,<address>,<age>,<religion>} in a text file by using a
        text editor. The file must be named: \texttt{raw\_input.txt}.
        An example text should be: \texttt{Davy Jones,2251234,168 Ly Thuong
        Kiet St District 10 Ward 14 HCMC,69,None}
  \item Open the file to read the text you inserted.
  \item Declare a string in the heap memory with dynamically allocated memory.
        The size of the string must be large enough to store the text in the
        text file.
  \item Copy the line from the text file to the string in the memory.
  \item Print the string to the terminal in the following format:
        \begin{enumerate}
          \item -{}-{}-{}-{}-Student personal information-{}-{}-{}-{}-
          \item Name: \textless name\textgreater
          \item ID: \textless ID\textgreater
          \item Address: \textless address\textgreater
          \item Age: \textless age\textgreater
          \item Religion: \textless religion\textgreater
        \end{enumerate}
        The output in the terminal for the former example input should be:\\
        \texttt{-{}-{}-{}-{}-Student personal information-{}-{}-{}-{}-}\\
        \texttt{Name: Davy Jones}\\
        \texttt{ID: 2251234}\\
        \texttt{Address: 168 Ly Thuong Kiet St District 10 Ward 14 HCMC}\\
        \texttt{Age: 69}\\
        \texttt{Religion: None}
  \item The printed result should also be written in a file named
        \texttt{formatted\_result.txt}
\end{enumerate}
\end{sloppypar}

\task{Test cases}
For each test case, \texttt{raw\_input.txt} is placed in the working directory of
MARS. After the run, the content of \texttt{formatted\_result.txt} is the same as
the terminal output.

\begin{sloppypar}
\testcase{3}{1}{\texttt{raw\_input.txt} contains\newline
  \texttt{Davy Jones,2251234,168 Ly Thuong Kiet St District 10 Ward 14 HCMC,69,None}}%
  {\newline
  \texttt{-{}-{}-{}-{}-Student personal information-{}-{}-{}-{}-}\newline
  \texttt{Name: Davy Jones}\newline
  \texttt{ID: 2251234}\newline
  \texttt{Address: 168 Ly Thuong Kiet St District 10 Ward 14 HCMC}\newline
  \texttt{Age: 69}\newline
  \texttt{Religion: None}}
\testcase{3}{2}{\texttt{raw\_input.txt} contains\newline
  \texttt{Nguyen Van An,2512345,268 Ly Thuong Kiet St Ward 14 District 10 HCMC,19,Buddhism}}%
  {\newline
  \texttt{-{}-{}-{}-{}-Student personal information-{}-{}-{}-{}-}\newline
  \texttt{Name: Nguyen Van An}\newline
  \texttt{ID: 2512345}\newline
  \texttt{Address: 268 Ly Thuong Kiet St Ward 14 District 10 HCMC}\newline
  \texttt{Age: 19}\newline
  \texttt{Religion: Buddhism}}
\testcase{3}{3}{\texttt{raw\_input.txt} contains (the line ends with a Windows line break \texttt{\textbackslash r\textbackslash n})\newline
  \texttt{Tran Thi Bich Ngoc,2310001,1 Vo Van Ngan St Linh Chieu Ward Thu Duc City,21,Christianity}}%
  {\newline
  \texttt{-{}-{}-{}-{}-Student personal information-{}-{}-{}-{}-}\newline
  \texttt{Name: Tran Thi Bich Ngoc}\newline
  \texttt{ID: 2310001}\newline
  \texttt{Address: 1 Vo Van Ngan St Linh Chieu Ward Thu Duc City}\newline
  \texttt{Age: 21}\newline
  \texttt{Religion: Christianity}}
\end{sloppypar}
